\documentclass[10pt,a4paper]{amsart}
\usepackage[margin=19mm]{geometry}
\usepackage{amsmath,amssymb,mathtools}
\usepackage[scr=rsfs,cal=euler]{mathalfa}
\usepackage{xfrac}
\usepackage[unicode,hidelinks,bookmarks=false]{hyperref}
\newcommand{\ie}{i.e.\ }
\newcommand{\cf}{cf.\ }
\newcommand{\resp}{respectively\ }
\newcommand{\Subsection}{\subsection}
\providecommand{\nobreakdash}{\nobreak\hbox{-}}
\setlength{\emergencystretch}{2em}
\multlinegap0pt
\usepackage{bm}
\usepackage{bbm}
\usepackage{dsfont}
\usepackage{xspace}
\usepackage{cancel}
\usepackage{mathtools}
\usepackage{amsthm}
\makeatletter
\@ifundefined{lem}{\newtheorem{lem}{Lemma}}{}
\@ifundefined{thm}{\newtheorem{thm}{Theorem}}{}
\makeatother
\newcommand{\rb}[1]{_{\rm #1}} % roman subscript 
\providecommand{\Nset}{\mathbb{N}}
\title[Combinations without specified separations]{Combinations without specified separations}
\author{Michael A. Allen}
\date{Source: arXiv:2210.08167v3 (CC BY 4.0)}
\begin{document}
\maketitle

\noindent\emph{Source and licence.} Combinations without specified separations. Version arXiv:2210.08167v3 (CC BY 4.0), distributed under the Creative Commons Attribution 4.0 International licence, as recorded in the arXiv licence metadata for this exact version and shown on the abstract page https://arxiv.org/abs/2210.08167v3. Full rights notice: licenses/arxiv-2210.08167-CC-BY-4.0.txt.\par\smallskip

\noindent\textit{Editorial scope.} Counted unit: the Thm whose source label is \texttt{T:CN} in the article (source0.tex lines 527--555, source label \texttt{T:CN}), delivered together with the article's own complete original proof of it (source0.tex lines 556--616, one \texttt{proof} environment of 61 lines). No mathematics is written, repaired or re-derived: statement and proof are the article's own text line for line. Required context retained uncounted: L:rr (lines 407--420). Every definition, display and label of those lines is retained unchanged. Omitted from this package and not redistributed: 1--406, 421--526, 617--1095. Nothing inside a retained excerpt is omitted or altered: every retained body is delivered exactly as the article's lines are written, so no omission receipt of any kind accompanies this package. Citations: the retained excerpts carry no citation command at all, so no bibliography entry of the article is transcribed and no reference is redistributed. Reference adaptation: none was needed and none was made; every reference inside the delivered unit resolves inside it, so no unresolved reference or citation is produced. Printed slips kept literal: no printed word, symbol, hypothesis or number of the counted statement or of its proof was altered. Layout and class: the article's own class is \texttt{elsarticle} with packages including graphicx, amsmath, amssymb, ifthen, color, hyperref; the wrapper is the lane's pinned \texttt{amsart} editorial wrapper at 10pt on A4, which restates the theorem-like environments the retained text needs and adds the article's own macro definitions that the retained text needs (\texttt{\textbackslash Nset}, \texttt{\textbackslash rb}), with extra wrapper packages (bm, bbm, dsfont, xspace, cancel, mathtools). The delivered numbering is the unit's own. Assets: no retained span contains a figure environment, a graphics-inclusion command, a PDF-page inclusion, a table or a TikZ picture; no image file is copied into this package, the package declares no asset dependency and no image file is redistributed with it. Licence: version arXiv:2210.08167v3 of this article is distributed under the Creative Commons Attribution 4.0 International licence, as recorded in the arXiv licence metadata for this exact version and shown on the abstract page https://arxiv.org/abs/2210.08167v3. A full rights notice is delivered at licenses/arxiv-2210.08167-CC-BY-4.0.txt. Characters: every retained excerpt is pure ASCII, so the delivered engine, XeLaTeX with its default Latin Modern text font, reports no missing character.
\begin{lem}\label{L:rr}
For all integers $n$ and $k$,
\begin{subequations}
\label{e:rr} 
\begin{align}\label{e:rrBn}
B_n&=\delta_{n,0}+\sum_{i=1}^{N\rb{m}}B_{n-l_i},\\
\label{e:rrBnk}
B_{n,k}&=\delta_{n,0}\delta_{k,0}+\sum_{i=1}^{N\rb{m}}B_{n-l_i,k-k_i},
\end{align}
\end{subequations}
where $N\rb{m}$ is the number of metatiles, $l_i$ is the length of the
$i$-th metatile and $k_i$ is the number of combs it contains, and
$B_{n<0}=B_{n,k<0}=B_{n<k,k}=0$.
\end{lem}

\begin{thm}\label{T:CN}
For a digraph possessing a common node, let $l_{\mathrm{o}i}$ be the
length of the $i$-th outer cycle ($i=1,\ldots,N\rb{o}$) and
let $k_{\mathrm{o}i}$ be the number of combs it contains, let $L_r$ be
the length of the $r$-th inner cycle ($r=1,\ldots,N$) and let $K_r$ be the
number of combs it contains, and let $l_{\mathrm{c}i}$ be the length
of the $i$-th common circuit ($i=1,\ldots,N\rb{c}$) and
let $k_{\mathrm{c}i}$ be the number of combs it contains. Then for all
integers $n$ and $k$,
\begin{subequations}
\label{e:CN}
\begin{align} 
\label{e:CNBn}
B_n&=\delta_{n,0}+ 
\sum_{r=1}^N (B_{n-L_r}-\delta_{n,L_r})+
\sum_{i=1}^{N\rb{o}}
\biggl(B_{n-l_{\mathrm{o}i}}-\sum_{r=1}^NB_{n-l_{\mathrm{o}i}-L_r}\biggr)
+\sum_{i=1}^{N\rb{c}} B_{n-l_{\mathrm{c}i}},\\
\label{e:CNBnk}
B_{n,k}&=\delta_{n,0}\delta_{k,0}+ 
\sum_{r=1}^N (B_{n-L_r,k-K_r}-\delta_{n,L_r}\delta_{k,K_r})+
\sum_{i=1}^{N\rb{o}}
\biggl(B_{n-l_{\mathrm{o}i},k-k_{\mathrm{o}i}}
-\sum_{r=1}^NB_{n-l_{\mathrm{o}i}-L_r,k-k_{\mathrm{o}i}-K_r}\biggr)\nonumber\\
&\qquad\mbox{}+\sum_{i=1}^{N\rb{c}} B_{n-l_{\mathrm{c}i},k-k_{\mathrm{c}i}},
\end{align} 
\end{subequations}
where $B_{n<0}=B_{n,k<0}=B_{n<k,k}=0$. 
\end{thm}

\begin{proof}
From Lemma~\ref{L:rr},
\begin{subequations}
\begin{align}\label{e:condCNn}
B_n&=\delta_{n,0}+\sum_{i=1}^{N\rb{o}}B_{n-l_{\mathrm{o}i}}+
\sum_{i=1}^{N\rb{c}}\sum_{j_1,\ldots,j_N\ge0}\!\!
\binom{j_1+\cdots+j_N}{j_1,\ldots,j_N}B_{n-\lambda_i},\\
\label{e:condCNnk}
B_{n,k}&=\delta_{n,0}\delta_{k,0}
+\sum_{i=1}^{N\rb{o}}B_{n-l_{\mathrm{o}i},k-k_{\mathrm{o}i}}+
\sum_{i=1}^{N\rb{c}}\sum_{j_1,\ldots,j_N\ge0}\!\!
\binom{j_1+\cdots+j_N}{j_1,\ldots,j_N}B_{n-\lambda_i,k-\kappa_i}
\end{align}
\end{subequations}
with $B_{n<0}=B_{n,k<0}=B_{n<k,k}=0$, 
where $\lambda_i=l_{\mathrm{c}i}+\sum_{s=1}^Nj_sL_s$ and 
$\kappa_i=k_{\mathrm{c}i}+\sum_{s=1}^Nj_sK_s$. 
The multinomial coefficient (which counts the number of
arrangements of the inner cycles) results from the fact that changing
the order in which the inner cycles are traversed (after the common
node is reached via the outgoing path of a common circuit) gives rise to
distinct metatiles of the same length.  The sum of terms over $j_1,\ldots,j_N$
in \eqref{e:condCNn} may be re-expressed as
\[
\sum_{j_1,\ldots,j_N\ge0}\!\!M(\varnothing)B_{n-\lambda_i}=
B_{n-l_{\mathrm{c}i}}+
\sum_{m=0}^{N-1}\sum_{\mathcal{R}_m}
\sum_{\substack{j_{s\notin\mathcal{R}_m}\geq1,\\j_{t\in\mathcal{R}_m}=0}}\!\!\!\!
M(\mathcal{R}_m)B_{n-\lambda_i}
\]
where $M(\mathcal{R})$ denotes the multinomial coefficient
$\tbinom{j_1+\cdots+j_N}{j_1,\ldots,j_N}$ with $j_{t\in\mathcal{R}}=0$,
and $\mathcal{R}_m$ denotes a set of $m$ numbers drawn from $\Nset_N$.  
For example, if $N>2$ an instance of $\mathcal{R}_2$ is $\{1,2\}$ in which case
$M(\mathcal{R}_2)=\tbinom{j_3+\cdots+j_N}{j_3,\ldots,j_N}$. 
Replacing $n$ by $n-L_r$ in
\eqref{e:condCNn} gives
\[
B_{n-L_r}=\delta_{n,L_r}+\sum_{i=1}^{N\rb{o}}B_{n-l_{\mathrm{o}i}-L_r}+
\sum_{i=1}^{N\rb{c}}\sum_{j_1,\ldots,j_N\ge0}\!\!
M(\varnothing)B_{n-\lambda_i-L_r}.
\]
After changing $j_r$ to $j_r-1$, 
the sum of terms over $j_1,\ldots,j_N$ may be re-expressed as
\[
\sum_{\substack{j_r\geq1,\\j_{t\neq r}\geq0}}\!\!\!\!
M_r(\varnothing)B_{n-\lambda_i}
=\sum_{m=0}^{N-1}\sum_{\mathcal{R}_m^{[r]}}
\sum_{\substack{j_{s\notin\mathcal{R}_m^{[r]}}\geq1,\\j_{t\in\mathcal{R}_m^{[r]}}=0}}\!\!\!\!
M_r(\mathcal{R}_m^{[r]})B_{n-\lambda_i},
\]
where $M_r(\mathcal{R})$ denotes the multinomial coefficient
$M(\mathcal{R})$ with $j_r$ replaced by $j_r-1$ (so, for example,
$M_3(\{1,2\})=\tbinom{j_3+\cdots+j_N-1}{j_3-1,\ldots,j_N}$) and
$\mathcal{R}_m^{[r]}$ is a set of $m$ numbers none of which equal $r$ 
drawn from $\Nset_N$.  After subtracting $\sum_{r=1}^N B_{n-L_r}$ from
\eqref{e:condCNn} and using the result for multinomial coefficients
that $M(\mathcal{R})=\sum_{r\notin R}M_r(\mathcal{R})$, we obtain
\eqref{e:CNBn}. Similarly, subtracting $\sum_{r=1}^N B_{n-L_r,k-K_r}$ from
\eqref{e:condCNnk} gives \eqref{e:CNBnk}.
\end{proof}

\end{document}
